\documentclass[10pt,a4paper]{amsart}
\usepackage[margin=19mm]{geometry}
\usepackage{amsmath,amssymb,mathtools}
\usepackage[scr=rsfs,cal=euler]{mathalfa}
\usepackage{xfrac}
\usepackage[unicode,hidelinks,bookmarks=false]{hyperref}
\newcommand{\ie}{i.e.\ }
\newcommand{\cf}{cf.\ }
\newcommand{\resp}{respectively\ }
\newcommand{\Subsection}{\subsection}
\providecommand{\nobreakdash}{\nobreak\hbox{-}}
\setlength{\emergencystretch}{2em}
\multlinegap0pt
\usepackage{amssymb}
\usepackage{mathtools}
\usepackage{tikz}
\newtheorem{corollary}{Corollary}
\title{Sums of $k$-bonacci Numbers}
\author{Harold R. Parks, Dean C. Wills}
\date{Source: arXiv:2208.01224v1 (CC BY 4.0)}
\begin{document}
\maketitle

\noindent\emph{Source and licence.} Sums of $k$-bonacci Numbers. Version arXiv:2208.01224v1 (CC BY 4.0), distributed by arXiv under the Creative Commons Attribution 4.0 International licence, http://creativecommons.org/licenses/by/4.0/, as recorded in the arXiv licence metadata for this exact version and shown on the abstract page https://arxiv.org/abs/2208.01224v1. Full licence notice: \texttt{licenses/arxiv-2208.01224-CC-BY-4.0.txt}.\par\smallskip

\noindent\textit{Editorial scope.} Counted unit: the article's corollary main.cor (source0.tex lines 567--585). The article's complete original proof of it is delivered with it (source0.tex lines 586--645, one \texttt{proof} environment of 60 lines). No mathematics is written, repaired or re-derived: the statement and the proof are the article's own lines, in the article's own notation, and no retained line is substituted or re-flowed.  Retained uncounted as the source-local context the proof needs: lines 134--141; lines 157--160; lines 542--547.  Nothing was omitted from any retained span, so the package carries no omission record.  No retained line contains a citation, so the wrapper carries no bibliography and no bibliography entry.  No retained line contains a figure environment, an image-inclusion command or drawing code, so no image is delivered and the package declares no asset dependency.  Editorial formatting only, in addition: an AMS \texttt{amsart} preamble that restates the article's own declarations and macros used by the retained lines, namely the environments \texttt{corollary} on separate counters, together with the standard packages those lines need.  The retained environments print in this document as Corollary 1 in order of appearance, whereas the article prints the same items with the numbering of its own full document.  The complete native source of the article as distributed in the arXiv e-print for this exact version is bundled unchanged as \texttt{sources/128979/source0.tex}.
\begin{equation}\label{fib.recur}
f_n^{(k)} = 
\begin{cases}
       0, &   n < 0, \\
       1, &   n = 0, \\
            \sum_{i=1}^k f_{n-i}^{(k)}, & n\geq 1.
\end{cases}
\end{equation}

\begin{equation}\label{main.comb}
\sum_{i=0}^{n} f_{i}^{(k)} = \sum_{j=0}^{\lfloor n/(k+1) \rfloor} (-1)^{j} \binom{n-jk}{j}2^{n-j(k+1)} 
\,.
\end{equation}

\begin{equation}
\label{eqn.comb2}
\sum_{i=0}^{n} f_i^{(k)}  %  S_{n+(k-1)}^{(k)}      
=\sum_{j=0}^{m} (-1)^{j} \binom{n-jk}{j}2^{n-j(k+1)} 
\,.
\end{equation}

\begin{corollary}\label{main.cor}
%  For  $k\geq 1$ and $n \geq 1$, we have
For  $k\geq 1$ and $n \geq 0$, we have\footnote{To avoid  $0/0$  when
$n=j=0$, we use the 
Kronecker $\delta$ defined by $\delta_{i, j} = 1$
if $i=j$ and $\delta_{i, j} = 0$ if $i\neq j$.}
\begin{equation}\label{cor.eqn1}
f_{n}^{(k)}= 
\sum_{j=0}^{\lfloor n/(k+1)\rfloor} (-1)^j\,
\frac{(n-jk)+j+\delta_{n,0}}{2(n-jk)+\delta_{n,0}}\,  \binom{n-jk}{j} \,  2^{n-j(k+1)}
\,.
\end{equation}
%\begin{equation}\label{cor.eqn}
%F_{n+(k-1)}^{(k)}= 
%\sum_{j=0}^{\lfloor n/(k+1)\rfloor} (-1)^j\,
%\frac{(n-jk)+j}{2(n-jk)}\,  \binom{n-jk}{j} \,  2^{n-j(k+1)}
%\,.
%\end{equation}
\end{corollary}

\begin{proof} 
For $n= 0$, the summation on the right-hand side of  (\ref{cor.eqn1})
consist of only the $j=0$ term, so  the result is confirmed by (\ref{fib.recur}).

For  $k\geq 1$ and $n \geq 1$,
%appropriate values of $n$ and $k$, 
using (\ref{main.comb}), we have:
\[
f_{n}^{(k)}=
\sum_{i=0}^{n} f_i^{(k)} 
- \sum_{i=1}^{n-1} f_i^{(k)}
%S_{n+(k-1)}^{(k)}-S_{(n-1)+(k-1)}^{(k)} 
\]
and we will use    (\ref{eqn.comb2}) 
to replace the partial sums of  $k$-bonacci numbers in this equation.
We will  need to examine the upper limits in the summations that
occur  in  (\ref{eqn.comb2}).

Set $q = \lfloor n/(k+1)\rfloor $.  Then $n = q(k+1) + r$, where
$r$ is an integer with $0\leq r < k+1$.
We  see that 
$$
\frac{n-1}{k} = q + \frac{q+r-1}{k}
\,.
$$
Since $n\geq 1$, one or both of $q$ and $r$ must be positive.
We conclude that $\lfloor (n-1)/k \rfloor \geq q$, so
%Note that $q\geq 1$.  Then  $\lfloor (n-1)/(k+1)\rfloor$ 
%may equal $q$ or $q-1$, but we have  
%$$
%n-1  \geq q(k+1) -1
%=qk +q -1\geq qk
%\,.
%$$
%Thus 
when    (\ref{eqn.comb2})  is applied with $n$ replaced by $n-1$
we may use $q$ as  the upper limit of summation.
%, because 
%$
%\lfloor (n-1)/k\rfloor  \geq q \geq   \lfloor (n-1)/(k+1)\rfloor
%$ holds.

We have
\begin{align}\nonumber
	f_{n}^{(k)} 
\nonumber
		&=\sum_{j=0}^{q} (-1)^{j} \binom{n-jk}{j}2^{n-j(k+1)} 
		- \sum_{j=0}^{q} (-1)^{j} \binom{(n-1)-jk}{j}2^{(n-1)-j(k+1)} \\
\nonumber %\label{use.pascal}
		&=\sum_{j=0}^{q} (-1)^{j}  \left[ 2 \binom{n-jk}{j}- \binom{n-jk-1}{j} \right] \frac12\, 2^{n-j(k+1)} \\
\nonumber
		&=\sum_{j=0}^{q} (-1)^{j} 
		\frac12 \left[ \frac{2(n-jk)}{n-jk}  - \frac{n-jk-j}{n-jk} \right] \binom{n-jk}{j} 2^{n-j(k+1)} \\
\nonumber
		&=\sum_{j=0}^{q} (-1)^{j} 
	 \frac{(n-jk) + j}{2(n-jk)}  \binom{n-jk}{j} 2^{n-j(k+1)} 
\end{align}
	 

\end{proof}

\end{document}
